Для практического применения систем искусственного интеллекта недостаточно обучить модель: необходимо организовать ее вывод, то есть выполнение модели на пользовательских запросах с заданными требованиями к задержке, пропускной способности и доступности. Такой этап работы модели принято называть инференсом. Современный распределенный инференс больших языковых моделей, основанных на архитектуре Mixture-of-Experts (MoE), представляет собой комплексную системную задачу. Многие существующие решения используют коллективные коммуникации и в первую очередь оптимизируют пропускную способность при фиксированном составе участников~\cite{aminabadi2022deepspeed,deepseekDeepEP}. Однако такая схема плохо подходит для динамической балансировки нагрузки между экспертными блоками: перераспределение запросов, использование реплик экспертов и добавление/исключение исполнителей на лету требуют более гибкой модели взаимодействия.

В рамках исследования предложена альтернативная клиент-серверная архитектура взаимодействия между основной частью модели и экспертными блоками. Разработан прототип системы, поддерживающий асинхронную диспетчеризацию запросов, мониторинг состояния экспертных блоков, динамический выбор доступного экспертного блока и обработку частичных отказов в процессе инференса. Для уменьшения накладных расходов, возникающих при замене коллективных операций адресным взаимодействием с экспертными блоками, реализован ряд оптимизаций, направленных на повышение эффективности передачи данных и вычислений на GPU.

Проведенные эксперименты показали, что предложенный подход обеспечивает возможность динамической балансировки нагрузки между экспертными блоками и сохранение доступности системы при частичных отказах за счет выбора доступных реплик экспертов. Полученные результаты подтверждают перспективность использования клиент-серверного подхода для построения гибких распределенных систем инференса MoE моделей~\cite{shazeer2017moe,deepseekDeepEP}.
